\documentclass[11pt]{article}
\usepackage{fontspec}
\setmainfont{Times New Roman}
\setsansfont{Arial}
\setmonofont{Consolas}
\usepackage{amsmath,amssymb,mathtools}
\usepackage{geometry}
\usepackage{microtype}
\geometry{a4paper,margin=2.35cm}
\setlength{\parindent}{1.25em}
\setlength{\parskip}{0pt}
\emergencystretch=25em
\sloppy
\hbadness=10000
\vbadness=10000
\hfuzz=2pt
\vfuzz=10pt
\raggedbottom

\begin{document}

% Унит: Папер29 в001. Дирецт Интерславиц Латин транслатион фром the Герман финал-аудитед соурце слице.
\section*{29. Теорем конечности инвариантов конечных линеарных груп карактеристики $p$}

\begin{center}
Nachr. в. d. Ges. d. Wiss. зу Готтинген 1926, стр. 28--35
\end{center}

Предложено од Р. Courant на заседаньу 30. јулија 1926.

В следујучем показујем, же знаны теорем конечности инвариантов конечных груп линеарных субституциј\footnote{За елементарны доказ конечности срав. Е. Noether, \emph{Der Endlichkeitssatz der Invarianten endlicher Gruppen}. Math. Ann. 77 (1916), стр. 89--92.} остаје важечим, ако се за област коефициентов, вместо обычного числового польа, положи в основу либовольно полье -- теды особливо полье карактеристики $p$.\footnote{За појмы теорије поль срав. Е. Steinitz, \emph{Algebraische Theorie der Körper}. J. f. M. 137 (1910), стр. 167--309. Карактеристика в § 4, кореньево полье в § 12, полье квоциентов в § 3.} Притом але знане доказы конечности одпадајут, скоро ред групы јест делимы чрез $p$; треба привлекти глубше аритметичне факты, имено следујучи критериј, доселе знаны само за карактеристику нульа:

\textbf{Критериј конечности}\footnote{Срав. Е. Noether, \emph{Algebraische und Differentialinvarianten}. Jahresber. d. d. Math.-Ver. 32 (1923), стр. 177--184, нр. 3, теорем 4 и там наведену литературу. В последньем ставку критерија тамо погрешно стоји „модулна база взходно к $\mathfrak R$” (фундаменталны систем целых величин) вместо „взходно к некому подколцу од $\mathfrak R$”.}: област целости $\mathfrak S$ из полиномов -- обчеје из рационалных функциј -- с коефициентами из польа јест конечна тогда и само тогда, когда в $\mathfrak S$ содержи се конечно подколцо $\mathfrak R$ тако, же вси елементи $\mathfrak S$ сут алгебраично целе над $\mathfrak R$. Всака интегрална база $\mathfrak R$ тогда дополньује се модулноју базоју $\mathfrak S$ взходно к некому подколцу од $\mathfrak R$ до интегралној базы $\mathfrak S$.

Тој критериј -- кторы сам по собе имаје самостојны интерес -- опираје се на ексистенцију рационалној базы и на факт, же теорем о цепах делительев остаје валидны при преходу к целым елементам конечного разширјеньа. Тај други теорем в претодној ноте Artin и van der Waerden доказали за либовольну карактеристику само под некојим предпоставјеньем конечности за изходну област, кторо јест исполньено в случају конечных груп (кореньево колцо представја конечно разширјенье). За ексистенцију рационалној базы дајем доказ валидны за либовольне польа како области коефициентов; с тым могу доказати критериј конечности под вышеј означеным ограниченьем.

Же за инварианты конечных груп критериј конечности јест исполньен, опираје се на принцип симетричных функциј. Но доколе при карактеристики нульа коефициенты Галоисовеј резолвенты уже дају интегралну базу, в обчем случају оны творет само конечно подколцо, потребно в критерију конечности. Из ексистенције того конечного подколца по тых самых закльучках следује конечност „релативных” инвариантов и „модуларных” инвариантов.

Меджу туто разгледане инварианты спадајут како специалне случаје „пројективне инварианты мод $p$”, обработане од Dickson и јего школы.\footnote{Срав. Dickson, \emph{On Invariants and the theory of Numbers}. The Madison Colloquium 1913. Новејша литература јест в томах Transactions of the American Mathematical Society, изданых одтого часа.} Тут был теорем конечности знаны за модуларне инварианты, але не за обче „формалне” инварианты.

\subsection*{§ 1. Критериј конечности}

1. \textbf{Теорем о ексистенцији рационалној базы.} Прво формулованье: Всакы систем $S$ рационалных функциј од $n$ неопредельеных, с коефициентами из польа $P$, имаје конечну рационалну базу; т. ј. обстаја конечно много функциј система тако, же всака функција система може се рационално, с коефициентами из $P$, изразити чрез те конечно многе функције.

Друго формулованье: Нехај $\mathfrak K$ буде меджуполье меджу $P$ и $P(x_1,\ldots,x_n)$ -- кде $x_1,\ldots,x_n$ означајут неопредельене -- трансцендентного степена $t\leq n$; нехај $y_1,\ldots,y_t$ буде неразложимы систем\footnote{Систем по Стеинитзу называје се неразложимы, ако он состоји из алгебраично независных функциј. Систем имаје трансцендентны степен $t$, ако содержи најмене једен неразложимы систем из $t$ елементов, але ниједен из више елементов. За просте теоремы о неразложимых системах, употребјене ниже, срав. Steinitz, § 22.} из $\mathfrak K$. Тогда $\mathfrak K$ јест конечно алгебраично разширјенье од $P(y_1,\ldots,y_t)$.

Оба формулованьа сут фактично еквивалентне. Достаточно јест доказати ексистенцију рационалној базы за вса меджупольа; бо полье выведено из система $S$ става тако меджуполье $\mathfrak K$, чија рационална база непосреднье даје рационалну базу од $S$, понеже кажды елемент $\mathfrak K$ става рационална комбинација конечно многих функциј из $S$. Из другого формулованьа следује пак, же либовольны неразложимы систем $y_1,\ldots,y_t$ из $\mathfrak K$, разширјен конечно многими функцијами, кторе карактеризујут $\mathfrak K$ како конечно разширјенье од $P(y_1,\ldots,y_t)$, даје таку рационалну базу. Обратно, ако по првом формулованьу јест дана рационална база од $\mathfrak K$, тогда по дефиницији трансцендентного степена кажды базовы елемент мора быти алгебраично зависны над $P(y_1,\ldots,y_t)$; с тым $\mathfrak K$ јест карактеризовано како конечно алгебраично разширјенье од $P(y_1,\ldots,y_t)$.

Теорем буде доказан в другом формулованьу. Најпрво нехај $\mathfrak K$ имаје тој самы трансцендентны степен $n$ како $P(x_1,\ldots,x_n)$; нехај $y_1,\ldots,y_n$ буде неразложимы систем из $\mathfrak K$ и следовательно из $P(x_1,\ldots,x_n)$. Због трансцендентного степена $n$ система $y_1,\ldots,y_n$ кажды $x_i$ става алгебраичны над $P(y_1,\ldots,y_n)$; зато $P(x_1,\ldots,x_n)$ и с тым $\mathfrak K$ како меджуполье става конечно алгебраично разширјенье од $P(y_1,\ldots,y_n)$.

Нехај тепер $t<n$; нехај нумерованье $x$ буде выбрано тако, же $y_1,\ldots,y_t,x_{t+1},\ldots,x_n$ творет неразложимы систем из $P(x_1,\ldots,x_n)$. Тогда $x_{t+1},\ldots,x_n$ сут неразложиме взходно к $P(y_1,\ldots,y_t)$ и зато по транзитивном закону алгебраичној зависности такоже неразложиме взходно к $\mathfrak K$ и к всакому меджупольу $\mathfrak M$ меджу $P(y_1,\ldots,y_t)$ и $\mathfrak K$. То значи, же $\overline{\mathfrak K}=\mathfrak K(x_{t+1},\ldots,x_n)$ односно $\overline{\mathfrak M}=\mathfrak M(x_{t+1},\ldots,x_n)$ става чисто трансцендентно разширјенье од $\mathfrak K$ односно $\mathfrak M$; зато кажды елемент из $\overline{\mathfrak K}$, кторы не принадлежи к $\mathfrak K$, става трансцендентны взходно к $\mathfrak K$, и подобно кажды елемент из $\overline{\mathfrak M}$, кторы не принадлежи к $\mathfrak M$, става трансцендентны взходно к $\mathfrak M$. Такоже нехај $\overline P$ буде положено равным чисто трансцендентному разширјеньу $P(x_{t+1},\ldots,x_n)$; тогда $y_1,\ldots,y_t$ сут неразложиме взходно к $\overline P$.

Доказ може се тепер привести к горнье решеному случају, ако $\overline P$ буде положено како област коефициентов. Из
\[
P<\mathfrak K<P(x_1,\ldots,x_t,x_{t+1},\ldots,x_n)
\]
\footnote{$P<\mathfrak K$ значи: $P$ содержи се в $\mathfrak K$, теды јест подполье од $\mathfrak K$.} следује
\[
\overline P<\overline{\mathfrak K}<\overline P(x_1,\ldots,x_t),
\]
при чем $\overline{\mathfrak K}$ имаје тој самы трансцендентны степен $t$ взходно к $\overline P$ како $\overline P(x_1,\ldots,x_t)$. Зато $\overline{\mathfrak K}$ взходно к $\overline P$ како области коефициентов имаје конечну рационалну базу, напр. $y_1,\ldots,y_t; z_1,\ldots,z_s$. Притом $z$ можно выбрати из $\mathfrak K$, понеже $\overline{\mathfrak K}$ јест полье породжено од $\mathfrak K$ и $\overline P$. Треба показати, же $\mathfrak K$ јест равно $\mathfrak M=P(y_1,\ldots,y_t;z_1,\ldots,z_s)$. $\mathfrak M$ јест меджуполье меджу $P(y_1,\ldots,y_t)$ и $\mathfrak K$; зато кажды елемент $\mathfrak K$ јест алгебраичны взходно к $\mathfrak M$. С друге страны, $\overline{\mathfrak M}$ јест равно $\overline{\mathfrak K}$ и зато $\mathfrak K$ содержи се в $\overline{\mathfrak M}$. Тогда $\mathfrak K$ јест меджуполье меджу $\mathfrak M$ и $\overline{\mathfrak M}$; а понеже кажды елемент из $\overline{\mathfrak M}$, кторы не принадлежи к $\mathfrak M$, јест трансцендентны взходно к $\mathfrak M$, нужно $\mathfrak K$ јест идентично с $\mathfrak M$.

Последок: Ако $P<\mathfrak K<\mathfrak L<P(x_1,\ldots,x_n)$ и $\mathfrak L$ јест алгебраично взходно к $\mathfrak K$, тогда $\mathfrak L$ јест конечно алгебраично разширјенье од $\mathfrak K$. Бо кажды елемент рационалној базы од $\mathfrak L$ тогда јест алгебраичны взходно к $\mathfrak K$; зато $\mathfrak L$ јест карактеризовано како конечно разширјенье од $\mathfrak K$.

2. \textbf{Доказ критерија конечности.} Условје јест јавно нужно; бо ако $\mathfrak S$ јест конечна област целости, тогда сама $\mathfrak S$ може се сматрјати како конечно подколцо $\mathfrak R$, наступајуче в критерију.

Условје јест достаточно под предпоставјеньем, же кореньево полье $P^{1/p}$ јест конечно разширјенье польа коефициентов\footnote{Гледај стр. 28, приметка 2.}, когда оно имаје карактеристику $p$; за карактеристику нульа никако ограничујуче предпоставјенье није потребно. За доказ треба показати, же $\mathfrak S$ имаје конечну модулну базу взходно к некому подколцу од $\mathfrak R$.

Нехај $\mathfrak K$ означа полье квоциентов\footnote{Гледај стр. 28, приметка 2.} од $\mathfrak R$, а $\mathfrak L$ полье квоциентов од $\mathfrak S$. Та польа квоциентов обстајут, понеже иде о колца без нуловых делительев; и имајемо $P<\mathfrak K<\mathfrak L<P(x_1,\ldots,x_n)$. Зато по последку под 1. $\mathfrak L$ става конечно алгебраично разширјено полье од $\mathfrak K$; далье, по предпоставјеньу, кажды елемент из $\mathfrak S$ јест целы взходно к $\mathfrak R$, и с тым $\mathfrak S$ става подколцо колца $\mathfrak S'$, состојајучего из всих елементов $\mathfrak L$, кторе сут целе над $\mathfrak R$. Но понеже ничто не јест предпоставјено о интегралној замкньености $\mathfrak R$ в $\mathfrak K$, факт, же теорем о цепах делительев остаје валидны при конечном разширјеньу, не може быти употребльен непосреднье. Најпрво треба прејдти к такоже конечному подколцу $\mathfrak T$ од $\mathfrak R$, за кторо та интегрална замкньеност јест исполньена.

Нехај најпрво област коефициентов $P$ содержи бесконечно много елементов. Тогда можно из $\mathfrak R$ выбрати неразложимы систем $g_1(x),\ldots,g_t(x)$ тако, же $\mathfrak R$, и с тым $\mathfrak S$, јест цела над $\mathfrak T=P[g_1(x),\ldots,g_t(x)]$. Имено, нехај $f_1(x),\ldots,f_k(x)$ буде интегрална база од $\mathfrak R$, ктора обстаја по предпоставјеньу; ако она не твори неразложимы систем, тогда -- по потребном пренумерованьу индексов -- нехај
\[
F(f_k;f_1,\ldots,f_{k-1})=0
\]
буде релација, ктореј $f_k$ задовольа. Ако заменити $f_1,\ldots,f_{k-1}$ чрез
\[
 f_1+t_1f_k,\ldots,f_{k-1}+t_{k-1}f_k,
\]
тогда $f_k$, и с тым такоже $f_1,\ldots,f_{k-1}$, јест целы над тими новыми функцијами, ако $t_i$ выбране сут како неопредельене; понеже $P$ имаје бесконечно много елементов, $t_i$ можно специализовати до елементов из $P$ тако, же та цела зависност остане. По конечно многих кроках, узимајучи в рачун транзитивны закон целој зависности, приходимо к исканому колцу $\mathfrak T=P[g_1(x),\ldots,g_t(x)]$. Тогда $\mathfrak L$ јест конечно разширјено полье польа квоциентов $P(g_1(x),\ldots,g_t(x))$ од $\mathfrak T$, и колцо $\mathfrak S'$ из всих елементов $\mathfrak L$, целых над $\mathfrak R$, јест идентично с колцом всих елементов $\mathfrak L$, целых над $\mathfrak T$.

В $\mathfrak T$ важи теорем о цепах делительев за систем всих идеалов, и $\mathfrak T$ јест интегрално замкньено в својем польу квоциентов; обоје зато, же $\mathfrak T$ јест изоморфно полиномној области од $t$ неопредельеных. Из предпоставјеньа, же $P^{1/p}$ јест конечно взходно к $P$, далье следује, же $\mathfrak T^{1/p}$ јест конечно взходно к $\mathfrak T$.\footnote{Срав. Hilbert, \emph{Über die vollen Invariantensysteme}, § 1. Math. Ann. 42 (1893), стр. 313--373.}\footnote{Срав. ноту Artin--van der Waerden цитовану в уводу.} Зато в $\mathfrak S'$ важи теорем о цепах делительев за систем всих $\mathfrak T$-модулов; особливо $\mathfrak S$, како подколцо од $\mathfrak S'$, обнимајуче $\mathfrak T$, имаје конечну $\mathfrak T$-модулну базу. Ексистенција конечној $\mathfrak T$-модулној базы за $\mathfrak S$ далье важи без никакого ограничујучего предпоставјеньа о области коефициентов $P$, ако $\mathfrak L$ става разширјенье првого рода од $P(g_1(x),\ldots,g_t(x))$, теды особливо всегда ако $P$ имаје карактеристику нульа.\footnote{Срав. напр. Е. Noether, \emph{Abstrakter Aufbau der Idealtheorie in algebraischen Zahl- und Funktionenkörpern}, § 3. Math. Ann. 96 (1926), стр. 26--61.}

Нехај $h_1(x),\ldots,h_s(x)$ буде $\mathfrak T$-модулна база од $\mathfrak S$. Тогда представјенье
\[
 f(x)=A_1(g)h_1(x)+\cdots+A_s(g)h_s(x)
\]
за либовольны $f(x)$ из $\mathfrak S$ -- кде $A$, како елементи из $\mathfrak T$, означајут полиномы в $g$ -- показује, же $g_1(x),\ldots,g_t(x),h_1(x),\ldots,h_s(x)$ творет искату интегралну базу.

Ако полье $P$ содержи само конечно много елементов, тогда полье $P(u)$, настало чрез адјункцију неопредельене $u$ -- т. ј. елемента трансцендентного взходно к $\mathfrak S$ -- содержи бесконечно много елементов. Зато колцо $\mathfrak S(u)$ имаје конечну интегралну базу взходно к $P(u)$ како области коефициентов; и понеже $\mathfrak S(u)$ јест колцо породжено од $\mathfrak S$ и $P(u)$, та база може быти выбрана из $\mathfrak S$ чрез раздельенье по степеньах $u$, при чем базове елементы $g_i$ од $\mathfrak S$ тепер не буду више творити неразложимы систем. С тым всакы полином $f(x)$ из $\mathfrak S$ допушча представјенье $C(u)f(x)=G(u,g_i,h_i)$, кде $C(u)$ означа ненуловы полином из $P[u]$ и $G$ става целы в $u,g_i,h_i$. Сравньенье коефициентов при входном степену $u$ показује, же в $g_i(x)$ и $h_i(x)$ дана јест интегрална база од $\mathfrak S$. С тым јест \textbf{критериј конечности доказаны в полној обчости}.

\subsection*{§ 2. Конечност инвариантов}

1. Под \textbf{конечноју линеарноју групоју карактеристики $p$} разумеје се група $\mathfrak H$ из $h$ линеарных субституциј $A_1,\ldots,A_h$ (с ненуловым детерминантом) с коефициентами из польа $P$ карактеристики $p$. Притом $A_k$ представја линеарну субституцију
\[
\begin{aligned}
 x_1^{(k)}&=a_{11}^{(k)}x_1+\cdots+a_{1n}^{(k)}x_n,\\[-0.2em]
 &\ldots,\\
 x_n^{(k)}&=a_{n1}^{(k)}x_1+\cdots+a_{nn}^{(k)}x_n,
\end{aligned}
\]
або скрачено: $(x^{(k)})=A_k(x)$. Група $\mathfrak H$ зато преводи ред $(x)$ с елементами $x_1,\ldots,x_n$ в реды $(x^{(k)})$ с елементами $x_1^{(k)},\ldots,x_n^{(k)}$. Понеже идентичност мора быти содержана меджу $A_1,\ldots,A_h$, меджу редами $(x^{(k)})$ содержи се такоже ред $(x)$.

Под \textbf{целым рационалным (абсолутным) инвариантом групы} разумеје се така цела рационална функција од $x_1,\ldots,x_n$ -- с коефициентами из $P$\footnote{То не огранича обчости, понеже всакы инвариант с коефициентами из польа карактеристики $p$ -- кторо мора имети обче надполье с $P$, абы инварианты были дефиниране -- може се линеарно разложити из конечно многих инвариантов с коефициентами из $P$. В обчем случају треба само изразити коефициенты чрез таке, кторе сут линеарно независне взходно к $P$; тогда својство инвариантности јест исполньено идентично в тых независных коефициентах.} -- ктора при применьеньу $A_1,\ldots,A_h$ остаје идентично неизменьена. К тим инвариантам зато принадлежат все симетричне функције -- с коефициентами из $P$ -- редов $(x^{(k)})$. Обратно, за всакы такы инвариант важи
\[
 f(x)=f(x^{(1)})=\cdots=f(x^{(h)})
 \quad\text{або}\quad
 h\cdot f(x)=f(x^{(1)})+\cdots+f(x^{(h)}).
\]
Из тога следује, же в случају, кде ред $h$ групы не јест делимы чрез карактеристику -- теды особливо в случају карактеристики нульа -- симетричне функције редов $(x^{(k)})$, с коефициентами из $P$, изчерпавајут целост инвариантов. Понеже те симетричне функције имајут конечну интегралну базу, доказ конечности в том случају јест доконаны; једночасно интегрална база добива се како систем $G_{\alpha\alpha_1\ldots\alpha_n}(x)$ коефициентов Галоисовеј резолвенты\footnote{Гледај стр. 28, приметка 1.}:
\[
 \Phi(z,u)=\prod_{k=1}^{h}\bigl(z-u_1x_1^{(k)}-\cdots-u_nx_n^{(k)}\bigr)
\]
\[
 =z^h+\sum G_{\alpha\alpha_1\ldots\alpha_n}(x)
 z^{\alpha}u_1^{\alpha_1}\cdots u_n^{\alpha_n}
 \quad(\alpha\ne h;\ \alpha+\alpha_1+\cdots+\alpha_n=h).
\]

2. В случају, же $h$ јест делимы чрез $p$, не јест нужно, же всакы инвариант групы јест симетрична функција редов $(x^{(k)})$, понеже тепер из $h\cdot f(x)$ не можно више закльучити на $f(x)$. Але колцо $\mathfrak R$, выведено из конечно многих инвариантов $G_{\alpha\alpha_1\ldots\alpha_n}(x)$ -- коефициентов Галоисовеј резолвенты -- твори конечно подколцо система $\mathfrak S$ всих инвариантов, и треба показати, же $\mathfrak S$ взходно к $\mathfrak R$ исполньа предпоставјеньа критерија конечности.

Најпрво без ограниченьа обчости\footnote{Гледај стр. 33, приметку.} за област коефициентов може се положити подполье $\overline P$ од $P$, выведено из всих коефициентов субституциј $a_{\mu\nu}^{(k)}$. То полье, выведено из конечно многих елементов, настаје чрез адјункцију конечно многих алгебраичных або трансцендентных елементов к првотному польу; и понеже првотно полье јест пополно, следује, же $\overline P^{1/p}$ јест конечно взходно к $\overline P$.\footnote{Гледај стр. 32, приметка 2.}

Далье треба показати, же $\mathfrak S$ јест алгебраично цела над $\mathfrak R$. По 1. меджу $h$ редами $x^{(k)}$ наступаје такоже изходны ред $x_1,\ldots,x_n$; зато Галоисова резолвента $\Phi(z,u)$ за $z=u_1x_1+\cdots+u_nx_n$ идентично изчезаје в $u$. Ако всакократ специализовати все $u$ к нули кроме $u_i$, кторо се поставльа равным јединици, добива се релација
\[
 x_i^h+\sum x_i^\alpha G_{\alpha0\ldots\alpha_i\ldots0}(x)=0
 \quad (\alpha\ne h;\ \alpha+\alpha_i=h),
\]
ктора показује, же кажды $x_i$ јест алгебраично целы над $\mathfrak R$. С тым исто важи за либовольны полином в $x$; особливо $\mathfrak S$ јест алгебраично цела над $\mathfrak R$. Тако сут предпоставјеньа критерија конечности исполньена, и \textbf{конечност инвариантов јест доказана}.

3. Треба приметити, же ты саме закльучки дају \textbf{конечност релативных и модуларных инвариантов}. Полином $f(x)$ называје се \textbf{релативны инвариант} групы, ако все $f(x^{(k)})$ разликују се од $f(x)$ само ненуловым фактором из $P$. Систем абсолутных инвариантов -- особливо колцо $\mathfrak R$, выведено из $G_{\alpha\alpha_1\ldots\alpha_n}(x)$ -- твори зато конечно подколцо колца выведеного из всих релативных инвариантов; и како горе, предпоставјеньа критерија конечности взходно к $\mathfrak R$ сут исполньена.

Полином $f(x)$ называје се \textbf{модуларны инвариант} групы, ако $f(x)$ става равно $f(x^{(k)})$ за всакы специалны систем вредностиј $x_1=\alpha_1,\ldots,x_n=\alpha_n$ из $P$. Всакы абсолутны инвариант јест једночасно модуларны инвариант; ако але $P$ содержи само конечно много елементов, обратно не мора важити. И тут сут предпоставјеньа критерија конечности исполньена взходно к подколцу $\mathfrak R$, выведеном из $G_{\alpha\alpha_1\ldots\alpha_n}(x)$.

\clearpage

\end{document}

